% Ported to the default house style by the handout-formatting skill, from
% inputs/math/Worksheets/02-Separable-Differential-Equations.pdf.
% One source, three audiences (change only this option line):
%   \usepackage[]{handout}        student handout (blanks, no solutions)
%   \usepackage[key]{handout}     answer key (solutions shown)
%   \usepackage[teacher]{handout} teacher copy (solutions + teacher notes)
\documentclass[10pt]{article}
\usepackage{housestyle}
\usepackage[teacher]{handout}

\hypersetup{pdftitle={Separable Differential Equations -- Worksheet 2}}

\begin{document}

\handouttitle{Separable Differential Equations}{Worksheet 2}{}

\calloutbox{Today's Goals}{%
    \begin{itemize}[leftmargin=*]
        \item Recognize when a first order equation \(\frac{dy}{dx}=H(x,y)\) is \textbf{separable}.
        \item Separate variables and integrate to find general and particular solutions.
        \item Use initial conditions and the sign of \(y\) to resolve absolute values and constants.
        \item Apply separable equations to exponential growth, radioactive decay, compound interest, and Newton's Law of Cooling.
    \end{itemize}%
}

\bigskip

A first order differential equation \(\dfrac{dy}{dx} = H(x,y)\) is called \textbf{separable}
if we can write \(H(x,y) = g(x)\,h(y)\). In that case we separate the variables and integrate:
\[
    \frac{dy}{dx} = g(x)\,h(y)
    \quad\Longrightarrow\quad
    \frac{1}{h(y)}\,dy = g(x)\,dx
    \quad\Longrightarrow\quad
    \int \frac{1}{h(y)}\,dy = \int g(x)\,dx.
\]

\section{Separating Variables}

\begin{enumerate}[leftmargin=*, label=\numbox{\arabic*}]

\item
    \begin{problem}[5cm]
        Solve the initial value problem \(\dfrac{dy}{dx} = -8xy\), \(\ y(0) = 4\).
    \end{problem}
    \begin{solution}
        Here \(H(x,y) = (-8x)(y) = g(x)h(y)\), so the equation is separable:
        \[
        \begin{aligned}
            \frac{dy}{y} &= -8x\,dx, &
            \int \frac{dy}{y} &= \int -8x\,dx, \\
            \ln|y| &= -4x^2 + c_3, &
            |y| &= e^{-4x^2}\cdot e^{c_3} = A e^{-4x^2}.
        \end{aligned}
        \]
        This is the general solution. Applying \(y(0) = 4\):
        \[
            4 = A e^{-4(0)^2} = A e^0 = A,
            \qquad\text{so}\qquad |y| = 4 e^{-4x^2}.
        \]
        Since \(y(0) = 4 > 0\), we have \(y > 0\) near \(x = 0\), so \(y = |y|\) and
        \[
            y = 4 e^{-4x^2} \qquad\text{(particular solution).}
        \]
    \end{solution}

\item
    \begin{problem}[6cm]
        Solve the initial value problem \(\dfrac{dy}{dx} = \dfrac{5 - 4x}{y(4y^2 + 2)}\), \(\ y(1) = 2\).
    \end{problem}
    \begin{solution}
        Writing the right side as \((5 - 4x)\dfrac{1}{4y^3 + 2y} = g(x)h(y)\), the equation is separable:
        \[
        \begin{aligned}
            (4y^3 + 2y)\,dy &= (5 - 4x)\,dx, \\
            \int (4y^3 + 2y)\,dy &= \int (5 - 4x)\,dx, \\
            y^4 + y^2 &= 5x - 2x^2 + c_3 \qquad\text{(general solution).}
        \end{aligned}
        \]
        This cannot easily be solved for \(y\) in terms of \(x\); it defines a family of curves.
        Applying \(y(1) = 2\):
        \[
            2^4 + 2^2 = 5(1) - 2(1)^2 + c_3
            \implies 20 = 3 + c_3 \implies c_3 = 17.
        \]
        So the solution is \(\ y^4 + y^2 = 5x - 2x^2 + 17\).
    \end{solution}

\item
    \begin{problem}[6cm]
        Find the general solution to \(\ y\,\dfrac{dy}{dx} - \left(8x^2 y\right)^{1/3} = 0\).
    \end{problem}
    \begin{solution}
        \[
        \begin{aligned}
            y\,\frac{dy}{dx} &= (8x^2 y)^{1/3} = 2x^{2/3} y^{1/3}, \\
            \frac{y}{y^{1/3}}\,\frac{dy}{dx} &= 2x^{2/3},
            \qquad y^{2/3}\,dy = 2x^{2/3}\,dx, \\
            \int y^{2/3}\,dy &= \int 2x^{2/3}\,dx, \\
            \frac{3}{5}y^{5/3} &= \frac{6}{5}x^{5/3} + c_3.
        \end{aligned}
        \]
        Multiplying by \(\tfrac{5}{3}\) gives \(y^{5/3} = 2x^{5/3} + c_4\), so
        \[
            y = \left(2x^{5/3} + c_4\right)^{3/5}.
        \]
    \end{solution}

\item
    \begin{problem}[6cm]
        Find the particular solution to the initial value problem
        \(\ e^{y}\dfrac{dy}{dx} = 3 e^{(3x - 2y)}\), \(\ y(0) = \tfrac{1}{3}\ln 5\).
    \end{problem}
    \begin{solution}
        \[
        \begin{aligned}
            e^{y}\frac{dy}{dx} &= 3 e^{(3x - 2y)} = \frac{3 e^{3x}}{e^{2y}}, &
            e^{3y}\frac{dy}{dx} &= 3 e^{3x}, \\
            \int e^{3y}\,dy &= \int 3 e^{3x}\,dx, &
            \frac{1}{3}e^{3y} &= e^{3x} + c_3.
        \end{aligned}
        \]
        Thus \(e^{3y} = 3 e^{3x} + c_4\), so \(3y = \ln(3 e^{3x} + c_4)\) and
        \[
            y = \frac{1}{3}\ln(3 e^{3x} + c_4) \qquad\text{(general solution).}
        \]
        Applying \(y(0) = \tfrac{1}{3}\ln 5\):
        \[
            \tfrac{1}{3}\ln 5 = \tfrac{1}{3}\ln(3 e^0 + c_4) = \tfrac{1}{3}\ln(3 + c_4)
            \implies 5 = 3 + c_4 \implies c_4 = 2.
        \]
        So \(\ y = \dfrac{1}{3}\ln(3 e^{3x} + 2)\) (particular solution).
    \end{solution}

\end{enumerate}

\section{Growth, Decay, and Cooling}

Both population growth and continuously compounded interest can be modeled by the rate of
change being a constant multiple of the amount: \(\dfrac{dP}{dt} = kP\), where \(k\) is the
growth rate. Separating variables gives \(\dfrac{1}{P}\,dP = k\,dt\), and since \(P(t) > 0\),
\[
    \ln P = kt + c_3 \implies P(t) = c_4 e^{kt} \implies P(t) = P_0 e^{kt},
\]
where \(P_0 = P(0)\). Radioactive decay obeys \(\dfrac{dN}{dt} = -kN\) with \(k \geq 0\), giving
\(N(t) = N_0 e^{-kt}\).

\begin{note}
    For continuous compounding, starting from \(A_0\) dollars at annual rate \(r\) compounded
    \(n\) times per year, the amount after \(t\) years is \(A_0\left(1 + \tfrac{r}{n}\right)^{nt}\).
    Letting \(n \to \infty\) gives \(A(t) = A_0 e^{rt}\), which satisfies \(\dfrac{dA}{dt} = rA\),
    \(A(0) = A_0\). Newton's Law of Cooling/Heating states that the rate of change of an object's
    temperature \(T\) immersed in a medium of constant temperature \(A\) is proportional to
    \(A - T\): \(\ \dfrac{dT}{dt} = k(A - T)\).
\end{note}

\begin{enumerate}[leftmargin=*, label=\numbox{\arabic*}, resume]

\item
    \begin{problem}[5cm]
        The population of a town in 2010 was \(100{,}000\). In 2013 its population was \(134{,}986\).
        (a) Find the annual growth rate. (b) How long does it take for the population to double?
    \end{problem}
    \begin{solution}
        \textbf{(a)} With \(P(t) = 100{,}000\, e^{kt}\) and \(P(3) = 134{,}986\):
        \[
            134{,}986 = 100{,}000\, e^{3k}
            \implies 1.34986 = e^{3k}
            \implies \ln(1.34986) = 3k \approx 0.3,
        \]
        so \(3k \approx 0.3\) and \(k = 0.1\). The annual growth rate is \(10\%\).

        \textbf{(b)} Doubling means \(200{,}000 = 100{,}000\, e^{0.1t}\):
        \[
            2 = e^{0.1t} \implies \ln 2 = 0.1t \implies t = 10\ln 2 \approx 6.93 \text{ years.}
        \]
    \end{solution}

\item
    \begin{problem}[5cm]
        Carbon~14 has an annual decay rate \(k = 0.0001216\). A piece of charcoal contains \(58\%\)
        as much Carbon~14 as a present-day sample of equal mass. What is the age of the sample?
    \end{problem}
    \begin{solution}
        Using \(N(t) = N_0 e^{-kt}\) with \(N(t) = 0.58 N_0\):
        \[
            0.58 N_0 = N_0 e^{-0.0001216 t}
            \implies \ln(0.58) = -0.0001216\, t
            \implies t = \frac{\ln(0.58)}{-0.0001216} \approx 4480 \text{ years old.}
        \]
    \end{solution}

\item
    \begin{problem}[5cm]
        Suppose you start with \(\$1000\) in an account compounded continuously at annual rate \(r\).
        After 3 years the amount has grown to \(\$1{,}116.28\). Find \(r\).
    \end{problem}
    \begin{solution}
        With \(A(t) = A_0 e^{rt}\), \(A_0 = 1000\), \(A(3) = 1116.28\):
        \[
            1116.28 = 1000\, e^{3r}
            \implies 1.11628 = e^{3r}
            \implies \ln(1.11628) = 3r,
        \]
        so \(r = \tfrac{1}{3}\ln(1.11628) \approx 0.0367\), i.e. \(r \approx 3.67\%\).
    \end{solution}

\item
    \begin{problem}[7cm]
        A roast, initially at \(40^\circ\)F, is placed in a \(400^\circ\)F oven at 5:00 pm. After
        90 minutes the roast is at \(150^\circ\)F.
        (a) Find a formula for the temperature after \(t\) minutes.
        (b) When will the roast reach \(160^\circ\)F?
    \end{problem}
    \begin{solution}
        \textbf{(a)} Solve \(\dfrac{dT}{dt} = k(400 - T)\). Separating variables (with \(400 - T > 0\)):
        \[
        \begin{aligned}
            \frac{1}{400 - T}\,dT &= k\,dt, &
            -\ln(400 - T) &= kt + c_3, \\
            \ln(400 - T) &= -kt - c_3, &
            400 - T &= c_4 e^{-kt},
        \end{aligned}
        \]
        so \(T(t) = 400 - c_4 e^{-kt}\). Using \(T(0) = 40\):
        \(\ 40 = 400 - c_4 \implies c_4 = 360\), giving \(T(t) = 400 - 360 e^{-kt}\).
        Using \(T(90) = 150\):
        \[
            150 = 400 - 360 e^{-90k}
            \implies \frac{25}{36} = e^{-90k}
            \implies k = -\frac{1}{90}\ln\!\left(\frac{25}{36}\right) \approx 0.00405.
        \]
        So \(T(t) = 400 - 360 e^{-0.00405 t}\).

        \textbf{(b)} Set \(T(t) = 160\):
        \[
            160 = 400 - 360 e^{-0.00405 t}
            \implies \frac{2}{3} = e^{-0.00405 t}
            \implies t = -\frac{1}{0.00405}\ln\!\left(\frac{2}{3}\right) \approx 100 \text{ minutes.}
        \]
        So the roast reaches \(160^\circ\)F at about 6:40 pm.
    \end{solution}

\end{enumerate}

\end{document}
